\chapter{Compatibilidade, Licença e Manutenção}

\section{Matriz de compatibilidade}

\begin{longtable}{p{0.27\textwidth}p{0.25\textwidth}p{0.38\textwidth}}
\toprule
\textbf{Documento ou componente} & \textbf{Versão} & \textbf{Situação} \\
\midrule
\endhead
Guia de Uso do UnDF-TeX & 1.0.0 & Compatível com a versão técnica indicada nesta edição. \\
Manual de TCC da UnDF & 1.0.0 & Documento acadêmico autônomo relacionado. \\
UnDF-TeX & 0.1.0 & Versão de referência deste guia. \\
LuaLaTeX e Biber & Distribuições atuais compatíveis & Devem oferecer os pacotes carregados pelo projeto. \\
\bottomrule
\end{longtable}

\section{Versionamento}

Mudanças na classe ou nos pacotes podem alterar comandos e aparência. Antes de atualizar um trabalho em andamento, registre a versão atual, crie uma cópia ou uma ramificação e compare os PDFs antes e depois. Não misture arquivos internos de versões diferentes.

Consulte o \href{https://github.com/igormrqs/undftex-publico/blob/main/CHANGELOG.md}{histórico público de alterações}
e a \href{https://github.com/igormrqs/undftex-publico}{página do repositório}
para identificar a origem do pacote. As versões editoriais 1.0.0 do manual
e do guia não representam, por si só, atualização do template para 1.0.0.

\section{Licença}

A LPPL 1.3c ou posterior aplica-se aos arquivos de código expressamente identificados. Consulte \texttt{LICENSE} e \texttt{LICENSE-LPPL-1.3c} antes de modificar ou redistribuir esses arquivos.

Marcas, logotipos, elementos de identidade institucional, fontes, imagens, documentos oficiais e recursos de terceiros permanecem sujeitos às respectivas regras de uso e não são automaticamente licenciados pela LPPL. A documentação textual somente deve ser redistribuída sob a licença que vier a ser expressamente indicada para ela.

\section{Responsabilidades de manutenção}

Uma alteração técnica deve atualizar código, exemplo, testes pertinentes, documentação de uso, histórico e indicação de compatibilidade. A validação deve incluir compilação do modelo em branco e do exemplo completo, comparação visual e revisão dos avisos do log.

\section{Relato de problema}

Informe versão, ambiente, entrada, comando, primeira mensagem e trecho mínimo. Em defeitos visuais, inclua página ou captura; omita dados pessoais e conteúdo sigiloso.

Problemas relativos à base pública podem ser registrados na
\href{https://github.com/igormrqs/undftex-publico/issues}{área de issues do repositório}.
Para recursos experimentais, use o canal de acompanhamento disponibilizado
pelo responsável pelo desenvolvimento privado.
